\documentclass[12pt]{article}

\usepackage[T1]{fontenc}
\usepackage[utf8]{inputenc}
\usepackage[brazil]{babel}
\usepackage{lmodern}
\usepackage{microtype}
\usepackage[a4paper,margin=2.5cm]{geometry}
\usepackage{setspace}
\usepackage{indentfirst}

\setstretch{1.15}

\title{Estatística e Governo \thanks{Texto Traduzido livremente por Luiz Mario Andrade com o auxílio da ferramenta Chat GPT 5.4}}
\author{Wesley C. Mitchell}
\date{Março de 1919}

\begin{document}

\maketitle

\begin{center}
    \textbf{ASSOCIAÇÃO AMERICANA DE ESTATÍSTICA} \\
    NOVA SÉRIE, No. 125. MARÇO, 1919
\end{center}

\vspace{1cm}

\begin{center}
    \textbf{ESTATÍSTICA E GOVERNO.*} \\
    Por Wesley C. Mitchell.
\end{center}

\vspace{0.5cm}

Desde que a Associação Americana de Estatística foi fundada em 1839, nenhum ano trouxe mudanças tão impactantes para a estatística americana quanto o ano que agora se encerra. A guerra forçou uma rápida expansão no escopo das estatísticas federais e a criação de novas agências estatísticas. O que é mais significativo, a guerra levou ao uso da estatística não apenas como um registro do que aconteceu, mas também como um fator vital no planejamento do que deveria ser feito. A guerra também trouxe um número sem precedentes de estatísticos para o emprego governamental. Provavelmente existem poucas sociedades profissionais que tiveram uma proporção tão considerável de seus membros engajados em trabalhos de guerra como esta Associação.

Esta noite, sentimos um justo orgulho pelo serviço que nossos colegas prestaram. Acalentamos a esperança de que aquilo que eles ajudaram a realizar durante a guerra em direção à orientação da política pública pelo conhecimento quantitativo do fato social não se perca no período de reconstrução pelo qual estamos passando, e no período indefinido de paz no qual estamos prestes a entrar. Para levar adiante essa esperança, a Associação pode buscar uma participação mais ativa no trabalho das estatísticas federais no futuro do que jamais teve no passado.

É minha agradável tarefa esta noite falar dessas coisas — nosso trabalho no ano que passou, nossas esperanças para os anos que virão e o que podemos fazer no presente para alcançar essas esperanças.

\footnotetext{* Discurso presidencial na octogésima Reunião Anual da Associação Americana de Estatística, Richmond, Virgínia, dezembro de 1918.}

\section*{I.}

A revisão do Dr. Cummings sobre as Estatísticas Federais em nosso Volume Comemorativo mostra que, antes da guerra, havia mais de vinte departamentos estatísticos em Washington publicando milhares de páginas de dados a cada ano. Mas não havia uma agência coordenadora, nenhum departamento cujo dever fosse considerar as necessidades estatísticas do governo como um todo e formular planos sistemáticos para suprir essas necessidades. Consequentemente, havia duplicação de trabalho, ninguém sabe exatamente quanta; muitos campos importantes eram imperfeitamente cobertos, ou nem sequer cobertos; os resultados de diferentes departamentos não podiam ser comparados ou combinados facilmente devido a diferenças nas unidades utilizadas, nos períodos cobertos e na classificação; finalmente, o custo das estatísticas federais era desnecessariamente alto. Para remediar esses defeitos de organização, Comitês de Estatísticas Federais nomeados por esta Associação e pela Associação Econômica esforçaram-se para garantir a nomeação de uma Comissão Estatística Central; mas encontraram pouco incentivo. Como membro de um desses comitês, confesso que nossos esforços não foram vigorosamente levados adiante.

A guerra revelou os defeitos da maquinaria federal para a coleta de estatísticas com uma rapidez surpreendente; pois a guerra impõe uma pressão sobre os escritórios estatísticos tanto quanto sobre as usinas de aço ou estaleiros. Como o Professor Young disse em seu discurso presidencial no ano passado: "A guerra tornou-se um conflito de massas direcionadas — de agregados. Homens, dinheiro, munições, alimentos, ferrovias, navegação, matérias-primas e produtos manufaturados em grande variedade são recrutados para o serviço da nação. Os problemas do controle e uso eficazes para fins de guerra desses variados recursos nacionais dependem intimamente do conhecimento de suas quantidades, isto é, da estatística... Assim como esta guerra é o nosso maior empreendimento nacional, suas demandas estatísticas constituem, no conjunto, o maior problema estatístico com o qual tivemos que lidar".

Não estávamos preparados para lidar com esse problema. Não se deve esperar, é claro, que a produção estatística de anos pacíficos inclua todos os dados necessários para travar uma guerra. Mas é de se esperar que uma organização governamental para a coleta de estatísticas compreenda um novo problema estatístico prontamente e prepare planos para tratar esse problema com vigor. Nesse teste, nossos departamentos federais falharam.

A falha foi enfaticamente uma falha de organização, em vez de funcionários individuais. Quaisquer acusações de incapacidade que sejam feitas contra os próprios funcionários devem ser devidamente feitas contra o sistema sob o qual os estatísticos federais são escolhidos e recompensados. Pois eles não são escolhidos com um olhar exclusivo para a habilidade técnica e capacidade administrativa; não recebem salários suficientes para atrair e reter homens de habilidade e ambição incomuns; os salários inadequados não são compensados pelo reconhecimento público de um serviço eficiente. Tínhamos, de fato, muitos estatísticos federais melhores do que o nosso tratamento merecia — homens que serviram ao país com zelo e inteligência. Mas, espalhados por numerosos pequenos departamentos, prescritos por uma rotina de deveres departamentais, com dotações escassas, esses homens tiveram pouca chance de considerar os vastos novos problemas da guerra. Eles certamente não apresentaram, e talvez não pudessem apresentar, um programa de guerra eficiente.

Por essa deficiência de nossa organização estatística, pagamos uma penalidade pesada. O tempo que gastamos na formulação de nossa organização de guerra e em colocá-la em funcionamento poderia ter sido substancialmente reduzido se alguém em Washington fosse capaz de apresentar prontamente às autoridades responsáveis os dados de que precisavam sobre homens e mercadorias, navios e fábricas.

O que aconteceu foi uma exibição admirável de energia e patriotismo nacional, mas não uma boa exibição de inteligência nacional. As Juntas de Guerra que o governo estabeleceu para suplementar os departamentos regulares enfrentaram tarefas estupendas. Eram lideradas e tripuladas, em sua maioria, por homens inexperientes em administração pública e desconhecedores dos deveres e recursos dos departamentos federais. Enquanto esses homens estavam em meio à elaboração de seus planos e formação de suas equipes, eles também tiveram que descobrir que precisavam de estatísticas, quais estatísticas precisavam e como obtê-las. Embora o governo federal tenha entrado na guerra com vinte ou mais agências estatísticas, o Conselho de Defesa Nacional, a Administração de Alimentos, a Administração de Combustíveis, a Junta de Navegação, a Junta de Comércio de Guerra, a Administração Ferroviária e a Junta de Indústrias de Guerra, mais cedo ou mais tarde, estabeleceram cada uma uma nova e independente agência estatística para atender às suas necessidades especiais. O Departamento de Guerra e o Departamento da Marinha seguiram o exemplo. E essas agências, assim como as Juntas de Guerra que as criaram, tiveram que ser tripuladas com pessoas inexperientes em trabalho governamental e pouco familiarizadas com Washington.

Embora eu tenha sido um dos recrutas brutos convocados para o trabalho de emergência do governo, não posso deixar de falar das finas qualidades que as novas equipes estatísticas demonstraram. Cada grupo estudou as necessidades particulares da junta à qual servia e lançou-se ardentemente na tarefa de coletar dados. Os novos homens trabalharam com intensidade apaixonada. Não se deixaram abater por obstáculos. Onde não conseguiam obter dados definitivos, não hesitavam em estimar. O lema adotado por um dos líderes expressava o espírito de todos: "Não pode ser feito? Mas aqui está".

No entanto, o trabalho estatístico das Juntas de Guerra como um todo mostrou precisamente o mesmo defeito que a organização dos antigos departamentos estatísticos, e mostrou essa falha em um grau agravado. Cada nova agência trabalhava por si mesma para uma junta separada. Consequentemente, houve muita duplicação de esforço e, ao mesmo tempo, muitos campos importantes permaneceram inexplorados; os resultados alcançados por diferentes agências não podiam ser facilmente comparados ou combinados; e o custo foi desnecessariamente alto. Além disso, a energia das novas agências estatísticas e a pressa com que trabalhavam ampliaram uma falha menor do antigo sistema para grandes proporções. Essas novas agências queriam obter seus dados fundamentais das fontes originais; então enviaram questionários a empresários em uma verdadeira inundação. Muitas plantas industriais receberam papéis elaborados que foram solicitados a preencher e devolver pelo próximo correio, às dezenas e centenas. Frequentemente, diferentes questionários cobriam quase o mesmo terreno e, geralmente, exigiam não pouca investigação dentro da planta para coletar os dados solicitados. Despesas consideráveis e séria irritação foram causadas em todo o país por essa óbvia falha de organização em Washington.

Esse mal dos questionários trouxe de volta uma enxurrada de reclamações, cujos ecos chegaram aos chefes responsáveis das Juntas de Guerra. A eficiência da mobilização econômica parecia ameaçada; isso era um assunto mais sério do que o desperdício de fundos públicos. Os homens que estavam mais cientes da falta de coordenação no trabalho estatístico agora tinham um forte argumento. Passos foram imediatamente tomados para remediar uma falha que era patente há uma geração ou mais em uma base de paz. O chefe da Divisão de Planejamento e Estatística da Junta de Navegação foi colocado encarregado do Departamento de Pesquisa da Junta de Comércio de Guerra e, em seguida, da Divisão de Planejamento e Estatística da Junta de Indústrias de Guerra. Assim, três das novas agências estatísticas foram colocadas sob uma única direção. Mais tarde, o mesmo homem tornou-se presidente do Comitê Estatístico do Departamento do Trabalho e, finalmente, foi autorizado a formar um Departamento Central de Planejamento e Estatística. O Departamento Central estabeleceu uma central de coordenação de atividades estatísticas, nomeou homens de contato para manter contato com o trabalho estatístico de todas as Juntas de Guerra e de alguns dos antigos departamentos, e começou a supervisionar a emissão de questionários. Quando o armistício foi assinado, estávamos em um bom caminho para desenvolver, pela primeira vez, uma organização sistemática de estatísticas federais.

Nas primeiras semanas após a interrupção dos combates, parecia que o que havia sido ganho na organização estatística poderia ser perdido quase de imediato. A rápida desmobilização das Juntas de Guerra ameaçava varrer com elas seus departamentos estatísticos, ou dispersar os novos departamentos estatísticos entre os antigos e nos deixar novamente em confusão estatística — produzindo números em abundância, mas sem um plano estatístico geral. Mas, em um momento crítico, o Presidente Wilson aprovou um plano pelo qual o Departamento Central de Planejamento e Estatística se tornou a única agência estatística a servir os delegados americanos na Mesa da Paz. Assim, o Departamento Central recebeu uma sobrevida por alguns meses. Resta saber se este departamento ou algum sucessor que desempenhe as mesmas funções centralizadoras será tornado permanente.

\section*{II.}

Ao falar em seguida de nossas esperanças para o futuro, estou falando meramente como um membro da Associação Americana de Estatística. No entanto, acredito que a maioria dos membros de nossa Associação crê que as ciências sociais em geral e as estatísticas sociais em particular têm um grande serviço a prestar ao governo e, através do governo, à humanidade.

O episódio na organização estatística que esbocei, o efeito da guerra sobre nossa atitude em relação ao uso de fatos para a orientação da política, liga o estágio atual da civilização ao passado selvagem do homem. Antropólogos passaram a reconhecer que as catástrofes desempenharam um papel de liderança no avanço da cultura. O selvagem e o bárbaro são criaturas tão conservadoras que nada menos que uma catástrofe pode abalá-los em seus hábitos estabelecidos, torná-los críticos dos velhos tabus, levá-los a usar sua inteligência livremente. Na ciência física e na técnica industrial, é verdade, nos emancipamos amplamente da dependência selvagem de catástrofes para o progresso. Pois nesses campos de atividade, desenvolvemos o hábito de criticar velhas formulações, de testar o que nossos pais aceitaram, de experimentar. Continuamos descartando o bom em favor do melhor, mesmo quando não estamos sob pressão. O resultado é uma taxa de avanço bastante constante — um avanço tão regular que contamos com ele ao traçar planos para o futuro. Hoje temos certeza de que, daqui a dez anos, nossas ideias científicas atuais e nossa maquinaria industrial atual estarão em boa parte antiquadas. Na ciência e na indústria, somos radicais — radicais que confiam em um método testado. Mas em questões de organização social, retemos uma grande parte do conservadorismo característico da mente selvagem. Uma grande catástrofe pode nos forçar por um pouco de tempo a levar a sério os problemas da mobilização social. Enquanto sob estresse, fazemos progressos rápidos. Mas quando o estresse passa, recuamos agradecidos para a nossa fé confortável no pensamento que foi feito por nossos pais.

Sei que há pessoas fervorosas que contestarão essas afirmações, pelo menos no presente. Elas confiam que a explosão de fervor patriótico trazida pela guerra nos levará triunfantes para a frente por uma geração. Contam com o autossacrifício generoso que todas as classes demonstraram, a fina disciplina que nossos soldados e trabalhadores de guerra mantiveram, para resolver os problemas da paz como resolveram o problema da guerra. Certamente nunca seremos novamente precisamente o que éramos antes da guerra. Mas tão certamente não permaneceremos o que fomos durante a guerra. Todos estamos sujeitos a reações emocionais e, como John Dewey apontou, o estado de espírito produzido pelo retorno da paz difere daquele produzido pelo eclodir da guerra tão amplamente quanto a própria paz difere da guerra. Não, não podemos depender de nenhum resquício da "psicologia de guerra" para organizar a democracia na paz.

O "reformador social" nós sempre teremos conosco, é verdade. Ou melhor, a maioria de nós somos "reformadores sociais" de algum tipo. E todos admiramos as qualidades que compõem os líderes da reforma social — profunda simpatia pelos oprimidos, coragem para enfrentar o ridículo, zelo ardente diante da indiferença, tato e energia na condução de cruzadas. Mas uma sucessão indefinida de campanhas para garantir esta, aquela e outra reforma específica é o que temos tido por muito, muito tempo. Muitas das reformas nas quais nossos avós, nossos pais e nossa juventude depositaram seus corações foram alcançadas. No entanto, a história do passado em questões de organização social não é uma história que gostaríamos que tivesse continuado por mil e um anos. A reforma por agitação ou luta de classes é uma maneira espasmódica de avançar, desconfortável e dispendiosa de energia. Não somos inteligentes o suficiente para conceber um método de progresso mais estável e seguro?

Com certeza, não poderíamos manter a organização social como ela é, mesmo que quiséssemos. Não estamos emergindo dos riscos da guerra para um mundo seguro. Pelo contrário, o mundo é um lugar muito perigoso para uma sociedade estruturada como a nossa, e eu, pessoalmente, fico feliz com isso. Os perigos são aumentados pelo nosso próprio progresso na indústria e na democracia. Não faz muito tempo, um físico inglês enfatizou o fato de que a cristandade moderna está consumindo em um ritmo cada vez maior a energia armazenada durante longas eras nos campos de carvão, e pintou o destino duvidoso da espécie humana como dependente da corrida entre a ciência e o átomo. Não chegou o momento de aplicar nossa inteligência para fazer o balanço dos recursos que a terra ainda possui e desenvolver métodos de utilização que protejam nosso futuro? Quanto ao progresso democrático, sabemos que homens que podem ler e votar tornam-se cidadãos inquietos em um estado onde seu trabalho não é interessante para eles e onde suas recompensas não satisfazem seu senso de justiça. E tal é o estado atual dos negócios com milhões de americanos agressivos. Pode-se contar com eles para mudar as coisas pelo tumulto se as coisas não forem mudadas pelo método.

Considerando-nos todos juntos como um povo em um grupo de povos poderosos, nossa primeira e principal preocupação é desenvolver alguma maneira de realizar os processos infinitamente complicados da indústria moderna e do intercâmbio dia após dia, apesar de todo o tédio e fadiga, e ainda assim manter-nos interessados em nosso trabalho e contentes com a divisão do produto. Essa é uma tarefa de suprema dificuldade — uma tarefa que exige experimentação inteligente e planejamento detalhado, em vez de agitação ou luta de classes. O que falta para alcançar esse fim, de fato, não é tanto boa vontade quanto conhecimento — acima de tudo, conhecimento do comportamento humano.

Nossa melhor esperança para o futuro reside na extensão para a organização social dos métodos que já empregamos em nossos campos de esforço mais progressistas. Na ciência e na indústria, como disse, não esperamos por catástrofes para forçar novos caminhos sobre nós. Não confiamos no poder propulsor de grandes emoções. Confiamos, e com sucesso, na análise quantitativa para apontar o caminho; e avançamos porque estamos constantemente melhorando e aplicando tal análise.

Embora eu pense que o desenvolvimento da ciência social ofereça mais esperança para resolver nossos problemas sociais do que qualquer outra linha de esforço, não afirmo que essas ciências em seu estado atual sejam muito úteis. Elas são imaturas, especulativas, cheias de controvérsias. Seus expoentes mais enérgicos ainda estão no estágio de desenvolver novos "pontos de vista", começando tudo de novo em um plano diferente, em vez de levar adiante a análise de seus predecessores. Em parte, as ciências sociais representam não o que é, mas o que seus autores pensam que deveria ser. Em suma, as ciências sociais ainda são infantis. Nem temos qualquer garantia certa de que elas algum dia crescerão para a robusta maturidade, não importa o cuidado que lhes dediquemos. Existem trilhas cegas de especulação nas quais gerações passadas mineraram industrialmente por eras com pouco ganho. Talvez as ciências sociais provem ser mais parecidas com a metafísica do que com a mecânica, mais parecidas com a teologia do que com a química. A raça pode sempre moldar seus destinos maiores por uma luta confusa na qual a força e a fraude, as boas intenções, o zelo ardente e a regra de ouro são fatores mais potentes do que a medição e o planejamento. Aqueles de nós que se preocupam com as ciências sociais, então, estão engajados em um empreendimento incerto; talvez não conquistemos grandes tesouros para a humanidade. Mas certamente é nossa tarefa trabalhar nessa pista com toda a inteligência e energia que possuímos até que sua riqueza ou esterilidade seja demonstrada.

As ciências sociais, no entanto, cobrem um campo imenso, e não é provável que encontremos fracasso ou sucesso em todas as suas partes. As partes onde o esforço parece mais promissor agora são as partes nas quais esta Associação está particularmente interessada. A medição é uma das características marcantes da ciência em geral, seja no campo da matéria inorgânica ou dos processos vitais. A estatística social, que se ocupa da medição dos fenômenos sociais, tem muitas das características progressistas das ciências físicas. Ela mostra um progresso direto no conhecimento do fato, na técnica de análise e no refinamento de resultados. É amigável à formulação matemática. É capaz de prever fenômenos de grupo. É objetiva. Um estatístico geralmente está ou certo ou errado, e seus sucessores podem demonstrar qual dos dois. Os estatísticos não estão continuamente começando sua ciência tudo de novo, desenvolvendo novos pontos de vista. Onde um investigador para, o próximo investigador começa com coleções maiores de dados, com extensões em novos campos ou com métodos de análise mais poderosos. Em todos esses aspectos, a posição e as perspectivas das estatísticas sociais assemelham-se mais à posição e às perspectivas das ciências naturais do que às das ciências sociais.

Acima de tudo, a estatística social, mesmo em seu estado atual, é diretamente aplicável em uma ampla gama na gestão de assuntos práticos, particularmente nos assuntos do governo. E esse valor prático da estatística é facilmente demonstrável até mesmo para um executivo ocupado. Uma vez obtida uma declaração quantitativa dos elementos cruciais no problema de um funcionário, apresente-a de forma concisa, ilumine as tabelas com um gráfico ou dois, encaderne o memorando em uma capa atraente amarrada com um laço caprichado, e é o homem excepcional que rejeitará sua ajuda. Daí em diante, seu problema não será fazer com que suas estatísticas sejam usadas, mas atender às chamadas contínuas por mais números, e evitar que seu convertido tome suas estimativas mais literalmente do que você mesmo as toma.

Podemos muito bem acalentar altas esperanças para o futuro imediato das estatísticas sociais. Ao contribuir para um conhecimento quantitativo dos fatos sociais, ao colocar esse conhecimento à disposição de funcionários responsáveis, estamos contribuindo com uma parte crucialmente importante para alcançar a tarefa mais grave que confronta a humanidade hoje — a tarefa de desenvolver um método pelo qual possamos fazer progressos cumulativos na organização social.

\section*{III.}

O que a Associação Americana de Estatística pode fazer para realizar essas esperanças? É claro que isso cabe à Associação decidir; mas atrevo-me a submeter certas recomendações ao julgamento da Associação.

Meu apelo é que a Associação busque desempenhar um papel mais ativo nos assuntos públicos do que desempenhou no passado. Estamos realizando nossa octogésima reunião anual — poucas sociedades científicas neste país são tão antigas. Durante todos esses anos fomos principalmente uma sociedade acadêmica, estimando nosso assunto particular, criticando aqueles que o negligenciam ou fazem mau uso dele, ocasionalmente oferecendo conselhos, resumindo experiências, mas não participando agressivamente no cotidiano da prática estatística. Todo esse tipo de trabalho tem sido útil. Certamente as condições em Washington e nas capitais estaduais tornaram a participação de estranhos nas estatísticas oficiais extremamente difícil. Mas as condições mudaram um pouco e, se fizermos nossa parte com vigor, elas podem mudar mais.

Duas mudanças parecem-me especialmente promissoras. Uma é a participação ativa que muitos membros da Associação tiveram no trabalho de guerra. Esses homens não perderão inteiramente o interesse pelas estatísticas federais quando deixarem Washington. Pelos próximos anos, pelo menos, teremos um corpo de trabalhadores que conhece bem as condições sob as quais as figuras do governo são compiladas e usadas. Esses homens ajudarão a tornar a Associação prática em qualquer conselho que ela venha a oferecer. Por causa deles, temos maior capacidade para realizar um trabalho útil agora do que jamais tivemos antes. A Associação pode ser mais útil porque sabe mais e se importa mais com o que os departamentos governamentais fazem.

A segunda mudança está na atitude dos funcionários de Washington em relação ao trabalho de estranhos. Assim como aqueles de nós que estiveram temporariamente no serviço governamental ganharam uma percepção simpática das dificuldades enfrentadas pelos departamentos estatísticos permanentes, os membros dos departamentos permanentes tornaram-se mais familiarizados com o ponto de vista dos estatísticos externos. Eles ouviram nossas críticas; por sua vez, criticaram muitas de nossas sugestões para melhorar sua organização e práticas. Como resultado, eles sabem como utilizar nossos serviços melhor do que faziam antes da guerra. E eles não estão, creio eu, indispostos a anular o divórcio entre o estatístico em exercício e o crítico acadêmico e entrar em um novo relacionamento de compreensão mútua e cooperação.

Um sintoma dessa nova atitude é tão gratificante que não posso deixar de chamar atenção especial para ele. O Secretário de Comércio pediu ao presidente da Associação Econômica Americana e ao presidente da Associação de Estatística que nomeiem, cada um, um comitê de três para aconselhar o Diretor do Censo em questões de princípio estatístico e na seleção de especialistas estatísticos. Esse arranjo, espera-se, não será uma formalidade; mas um plano de trabalho pelo qual os produtores e os consumidores de estatísticas possam cooperar efetivamente para melhorar os produtos nos quais ambas as partes estão interessadas. Para fornecer aos dois comitês instalações de trabalho, um escritório e um secretário foram fornecidos a eles pelo Diretor do Censo.

Se fizermos nossa parte para tornar esse arranjo um sucesso, ele poderá talvez levar ao estabelecimento de outros vínculos entre as Associações que representam o público estatístico e os escritórios nos quais as estatísticas são preparadas.

Existem várias medidas práticas para as quais podemos contribuir, se quisermos. Por exemplo, podemos usar nossa influência sempre que surgir oportunidade para garantir salários mais adequados para os estatísticos do governo. A escala de pagamento era muito baixa antes da guerra; o aumento do custo de vida tornou-a chocantemente inadequada. A menos que aumentos sejam concedidos, muitos homens experientes que ficariam felizes em continuar no serviço público serão forçados, em justiça às suas famílias, a procurar vagas em outros lugares. Agora que a guerra acabou, não podemos justamente pedir que esses homens economizem em seus filhos pelo resto de nós. A profissão do estatístico exige habilidade e treinamento não menores do que os necessários para a contabilidade; no entanto, pelo que posso aprender, a remuneração média dos estatísticos é decididamente inferior à dos contadores. Como representante da profissão estatística, é certamente direito desta Associação instar vigorosamente uma escala salarial mais alta.

Podemos também tomar uma posição definida sobre a continuação das novas atividades estatísticas iniciadas durante a guerra. As Juntas de Guerra acharam necessário obter dados mensais de estoques de certas mercadorias em mãos e dados mensais da produção de outras mercadorias. Esses dados foram coletados de várias maneiras, pelo Departamento do Censo, por organizações comerciais como o Conselho de Curtidores, ou por seções das próprias Juntas de Guerra. Os resultados são de interesse não apenas para as indústrias envolvidas, mas também para o governo e para o público em geral. A manutenção permanente desse serviço, talvez de forma modificada, é uma medida que promete comandar o apoio crescente dos empresários. Se sistematicamente estendido, esse trabalho poderia muito bem se transformar em um censo contínuo de produção, simples na forma, barato; mas de grande valor na previsão das condições de negócios e na orientação da política pública.

Mais uma vez, há a questão que mencionei na primeira seção deste discurso, se o Departamento Central de Planejamento e Estatística deve ser continuado ou dissolvido quando a Conferência de Paz terminar seu trabalho. Alguma agência centralizadora para considerar as necessidades estatísticas do governo como um todo e para estabelecer planos sistemáticos para atender a essas necessidades é nossa maior lacuna estatística. Em uma questão desse caráter, não é dever da Associação Americana de Estatística expressar sua opinião?

Em qualquer ação que tomarmos, faremos bem em distinguir claramente entre dois tipos de estatísticas — as estatísticas que são usadas como um registro do que foi e as estatísticas que são usadas como base para o planejamento do que será. Desses dois tipos, as estatísticas de registro são as mais familiares. Elas constituem os números que entram em relatórios anuais, que são analisados minuciosamente pelo estudante, que são citados muito tempo depois pelo historiador. Tais figuras têm uma influência na moldagem de políticas públicas; mas essa influência é vaga e intermitente. O administrador médio pouco se importa com o que aconteceu antes de ontem — seus pensamentos estão obcecados pelo que está acontecendo hoje e o que deveria acontecer amanhã. Qualquer um de nós em sua posição desenvolveria esse estado de espírito se tivesse algum sucesso. O que o administrador precisa para orientar a política pública, o que ele aprenderá rapidamente a usar se os obtiver, são estatísticas de planejamento bem organizadas. As estatísticas de planejamento para serem úteis devem ser estritamente atualizadas. Elas devem mostrar os fatores vitais da situação. Devem ser apresentadas de forma concisa, em formato padronizado, tanto em gráficos quanto em tabelas. Os dados devem ser simples o suficiente para serem enviados por telégrafo e compilados da noite para o dia. Aproximações brutas servirão ao propósito. O que precisamos praticamente no presente é desenvolver agências estatísticas para obter tais estatísticas de planejamento e colocá-las diante dos homens cujas decisões são importantes para o país, sejam esses homens administradores, legisladores ou eleitores.

Como estudantes, nossa preocupação continuará a ser principalmente com as estatísticas de registro — elas não devem ser negligenciadas; de fato, devem ser estendidas e melhoradas. Mas como homens interessados nos assuntos públicos, nossa ênfase deve ser colocada no desenvolvimento e no uso de estatísticas de planejamento.

A política de participação ativa na moldagem do trabalho estatístico que estou urgindo parece-me justificada pelas circunstâncias atuais. Durante a guerra aprendemos que muitas coisas que pareciam impossíveis eram fáceis de realizar se atacadas com vigor. Sem dúvida, a situação já se cristalizou em parte; mas muitos assuntos da política governamental ainda estão em um estado fluido. Algumas mudanças terão que ocorrer; a questão é quais serão essas mudanças. Se colocarmos nosso conhecimento técnico e nossa experiência prática à disposição da nação, poderemos aumentar nossa influência pelos anos que virão e, o que é vastamente mais importante, poderemos ajudar a tornar o conhecimento quantitativo dos fatos um fator potente no governo.

\end{document}